\documentclass{article}%
\usepackage[T1]{fontenc}%
\usepackage[utf8]{inputenc}%
\usepackage{lmodern}%
\usepackage{textcomp}%
\usepackage{lastpage}%
%
\usepackage[russian]{babel}%
\usepackage{fontenc}%
\usepackage{graphicx}%
\usepackage{float}%
\usepackage{longtable}%
\usepackage{array}%
%
\begin{document}%
\normalsize%
\section{Тестовый документ}%
\begin{flushleft}%
Это обычный текст.%Шрифт-Times New Roman Размер шрифта-12.0 Межстрочный интервал-1.5 Правило-ONE_POINT_FIVE (1)%%
\end{flushleft}%
\begin{flushleft}%
\textbf{\textit{Это}}%Шрифт-Bahnschrift Light Размер шрифта-18.0 Межстрочный интервал-1.5 Правило-ONE_POINT_FIVE (1)% %Шрифт-Bahnschrift Light Размер шрифта-18.0 Межстрочный интервал-1.5 Правило-ONE_POINT_FIVE (1)%\textbf{\textit{жирный}}%Шрифт-Bahnschrift Light Размер шрифта-18.0 Межстрочный интервал-1.5 Правило-ONE_POINT_FIVE (1)% %Шрифт-Bahnschrift Light Размер шрифта-18.0 Межстрочный интервал-1.5 Правило-ONE_POINT_FIVE (1)%\textbf{\textit{и}}%Шрифт-Bahnschrift Light Размер шрифта-18.0 Межстрочный интервал-1.5 Правило-ONE_POINT_FIVE (1)% %Шрифт-Bahnschrift Light Размер шрифта-18.0 Межстрочный интервал-1.5 Правило-ONE_POINT_FIVE (1)%\textbf{\textit{курсивны}}%Шрифт-Bahnschrift Light Размер шрифта-18.0 Межстрочный интервал-1.5 Правило-ONE_POINT_FIVE (1)%\textbf{\textit{й}}%Шрифт-Bahnschrift Light Размер шрифта-18.0 Межстрочный интервал-1.5 Правило-ONE_POINT_FIVE (1)% %Шрифт-Bahnschrift Light Размер шрифта-18.0 Межстрочный интервал-1.5 Правило-ONE_POINT_FIVE (1)%\textbf{\textit{текст}}%Шрифт-Bahnschrift Light Размер шрифта-18.0 Межстрочный интервал-1.5 Правило-ONE_POINT_FIVE (1)%\textbf{\textit{.}}%Шрифт-Bahnschrift Light Размер шрифта-18.0 Межстрочный интервал-1.5 Правило-ONE_POINT_FIVE (1)%%
\end{flushleft}%
\begin{flushleft}%
\textbf{Это жирный текст.}%Шрифт-Arial Black Размер шрифта-8.0 Межстрочный интервал-16.0 Правило-EXACTLY (4)%%
\end{flushleft}%
\begin{flushleft}%
\textit{Это курсивный текст.}%Шрифт-None Размер шрифта-10.0 Межстрочный интервал-16.0 Правило-EXACTLY (4)%%
\end{flushleft}%
\begin{flushleft}%
Это список.%Шрифт-None Размер шрифта-None Межстрочный интервал-None Правило-None%%
\end{flushleft}%
\begin{enumerate}%
\item Это %Шрифт-None Размер шрифта-None Межстрочный интервал-None Правило-None%список%Шрифт-None Размер шрифта-None Межстрочный интервал-None Правило-None%.%Шрифт-None Размер шрифта-None Межстрочный интервал-None Правило-None%%
\item \textbf{\textit{Это список жирный и курсивный.}}%Шрифт-None Размер шрифта-None Межстрочный интервал-None Правило-None%%
\item \textbf{Это список жирный.}%Шрифт-None Размер шрифта-None Межстрочный интервал-None Правило-None%%
\item \textit{Это}%Шрифт-None Размер шрифта-None Межстрочный интервал-None Правило-None%\textit{ список курсивный.}%Шрифт-None Размер шрифта-None Межстрочный интервал-None Правило-None%%
\end{enumerate}%
\begin{longtable}{|>{\raggedright\arraybackslash}p{3.190875cm}|>{\raggedright\arraybackslash}p{8.251472222222223cm}|>{\raggedright\arraybackslash}p{5.207cm}|}%
\hline%
ячейка%Шрифт-None Размер шрифта-None Межстрочный интервал-None Правило-None% & \textbf{\textit{жирная и курсивная ячейка}}%Шрифт-None Размер шрифта-None Межстрочный интервал-None Правило-None% & \textbf{Жирная ячейка}%Шрифт-None Размер шрифта-None Межстрочный интервал-None Правило-None% \\ \hline \endfirsthead%
\hline%
\textit{Курсивная ячейка}%Шрифт-None Размер шрифта-None Межстрочный интервал-None Правило-None% & ячейка%Шрифт-None Размер шрифта-None Межстрочный интервал-None Правило-None% & ячейка%Шрифт-None Размер шрифта-None Межстрочный интервал-None Правило-None% \\ \hline%
ячейка%Шрифт-None Размер шрифта-None Межстрочный интервал-None Правило-None% & ячейка%Шрифт-None Размер шрифта-None Межстрочный интервал-None Правило-None% & ячейка%Шрифт-None Размер шрифта-None Межстрочный интервал-None Правило-None% \\ \hline%
\end{longtable}%
\begin{flushleft}%
Это текст по левому краю%Шрифт-None Размер шрифта-None Межстрочный интервал-None Правило-None%%
\end{flushleft}%
\begin{flushright}%
Это текст по правому краю%Шрифт-None Размер шрифта-None Межстрочный интервал-None Правило-None%%
\end{flushright}%
\begin{flushleft}%
Это текст по ширине%Шрифт-None Размер шрифта-None Межстрочный интервал-None Правило-None%%
\end{flushleft}%
\begin{center}%
Это текст по центру%Шрифт-None Размер шрифта-None Межстрочный интервал-None Правило-None%%
\end{center}%
\begin{figure}[H]%
\centering%
\includegraphics[width=0.8\textwidth]{images/img\_6.png}%
\end{figure}%
\section{Это заголовок 1 уровня}%
\subsection{Это заголовок 2 уровня}%
\subsubsection{Это заголовок 3 уровня}%
\paragraph{Это заголовок 4 уровня}%
\paragraph{Это заголовок 5 уровня}%
\end{document}